\documentclass[a4paper,12pt]{article}

\usepackage[utf8]{inputenc}
\usepackage{amsmath,amssymb}
\usepackage{geometry}
\geometry{margin=2.5cm}

\newenvironment{exercise}
  {\par\vspace{1em}\noindent\textbf{Esercizio.} }
  {\par\vspace{1em}}

\newenvironment{proof}
  {\par\vspace{0.5em}\noindent\textit{Dimostrazione.} }
  {\par\vspace{1em}\hfill$\square$}


\begin{document}

\begin{exercise}

	Suppose $B$ b.a. and denote $\tau$ the topology of $St(B)$. Consider
	\begin{itemize}
		\item	$I: \tau \to \{\text{ideals of }B\}, Z \mapsto \{b \in B : N_b \subseteq Z\} $ and $X: \{\text{ideals of }B \} \to \tau, I \mapsto \bigcup\{N_b : b \in I\}$.
			Then $I,X$ are one the inverse of the others, so there is a correspondence between open subsets of $St(B)$ and ideals of $B$.
			Furthermore $U \subseteq St(B)$ is open if and only if $X(I(U)) = U$.

		\item	$F: \tau \to \{\text{filters of }B\}, Z \mapsto \{b \in B : N_b \supseteq Z\}$ and 
			$X: \{\text{filters of }B\} \to \{\text{ closed subsets of St(B) }\}, F \mapsto \bigcap\{N_b : b \in F\}$.
			Then $F,X$ are one the inverse of the other, so there is a correspondence between closed subsets of $St(B)$ and filters of $B$.
			Furthermore $C \subseteq St(B)$ is closed if and only if $X(F(C)) = C$.

	\end{itemize}

\end{exercise}

\begin{proof}

	\begin{itemize}
		\item	Let $Z \subseteq St(B)$. Then $X(I(Z)) = X( \{b \in B : N_b \subseteq Z\} ) = \bigcup\{ N_b : N_b \subseteq Z\} = Z$, where the last inequality holds because $Z$ is open.
			Let $J$ ideal of $B$. Then $I(X(J)) = I( \bigcup\{N_b : b \in J\}) = \{c \in B : N_c \subseteq \bigcup\{N_b : b \in J\}\}$.
			Then it is clear that $I(X(J)) \supseteq J$. For the converse, let $c \in I(X(J))$. Then $N_c \subseteq \bigcup\{N_b : b \in J\}$.
			Since we know $St(B)$ is compact and $N_c$ is closed, we have $N_c$ compact, so we can extract a finite covering: $N_c \subseteq \bigcup_{i=1,..,k} N_{b_i} = N_{\bigvee_{i=1,..,k} b_i}$.
			So we deduce $c \leq \bigvee_{i=1,..,k} b_i \in J$ and since $J$ ideal we have $c \in J$ as well.
			This way we have incidentally proved that if $U \subseteq St(B)$ is open then $X(I(U)) = U$.
			For the converse observe that $X(I(U))$ is union of open subsets of $St(B)$.


		\item	Pretty much the same.

	\end{itemize}

\end{proof}
\end{document}
